\subsection{转置}

设 E 是局部凸空间，\(E^{\prime}\) 是它的对偶。回顾，我们用 \(E_{b}^{\prime}\) 和 \(E_{c}^{\prime}\) 分别表示给向量空间 \(E^{\prime}\) 赋以有界收敛拓扑和紧收敛拓扑所得到的局部凸空间（EVT, III, p. 14）。空间 \(E_{b}^{\prime}\) 也称为 E 的强对偶（同前引处, 例 4）。

命题 9. --- 设 E 是局部凸空间，F 是分离的局部凸空间，u 是从 E 到 F 的连续线性映射。

\begin{enumerate}
\def\labelenumi{\alph{enumi})}
\item
  若映射 \(u\) 是紧的，则它的转置 \(^{t}u\) 是从 \(F_c'\) 到 \(E_c'\) 的紧线性映射，且是从 \(F_c'\) 到 \(E_b'\) 的连续线性映射；
\item
  假设空间 E 是半赋范空间且空间 F 是拟完备的。若转置 \({}^{t}u\) 是从 \(F_{c}^{\prime}\) 到 \(E_{b}^{\prime}\) 的连续映射，则线性映射 u 是紧的。
\end{enumerate}

首先假设线性映射 \(u\) 是紧的。于是存在 E 中 0 的一个邻域 V，它在 \(u\) 下的像包含于 F 的一个紧子集 C。按照定义（EVT, II, p. 47, def. 2 及 p. 68, def. 1），C 的极集 \(\mathrm{C}^{\circ}\) 是 \(\mathrm{F}_c^\prime\) 中 0 的一个邻域，并且 \({}^t u(\mathrm{C}^\circ) \subset \mathrm{V}^\circ\)。而 \(\mathrm{V}^\circ\) 是 \(\mathrm{E}'\) 的一个等度连续且对弱拓扑为闭的子集，从而在 \(\mathrm{E}_c^\prime\) 中紧（EVT, III, p. 17, cor. 2 及 prop. 5）。这就证明了 \({}^t u\) 是从 \(\mathrm{F}_c^\prime\) 到 \(\mathrm{E}_c^\prime\) 的紧线性映射。

我们保留前面的假设，设 U 是 \(E_{b}^{\prime}\) 中 0 的一个邻域。它包含 E 的某个有界子集 A 的极集 \(A^{\circ}\)。由于线性映射 u 紧，集合 \(u(\mathrm{A})\) 在 F 中相对紧（III, p. 2, 注记 1），且 \((^{t}u)^{-1}(\mathrm{A}^{\circ}) = u(\mathrm{A})^{\circ}\) 是 \(F_{c}^{\prime}\) 中 0 的一个邻域。这就证明了 \(^{t}u\) 是从 \(F_{c}^{\prime}\) 到 \(E_{b}^{\prime}\) 的连续线性映射。

现在置于 \(b\) 的假设之下，并用 B 表示 E 的单位球。B 的极集 \(\mathrm{B}^{\circ}\) 是 \(\mathrm{E}_b'\) 的单位球，并且 \((^{t}u)^{-1}(\mathrm{B}^{\circ}) = u(\mathrm{B})^{\circ}\)。若 \(^{t}u\) 是从 \(\mathrm{F}_c'\) 到 \(\mathrm{E}_b'\) 的连续映射，则集合 \(u(\mathrm{B})^{\circ}\) 是 \(\mathrm{F}_c'\) 中 0 的一个邻域；于是它包含 F 的某个紧子集 C 的极集 \(\mathrm{C}^{\circ}\)。由双极定理（EVT, II, p. 49, cor. 3），集合 \(u(\mathrm{B})\) 包含于 \(\mathrm{C} \cup \{0\}\) 的闭凸包；由于空间 F 拟完备（EVT, III, p. 8），后者是紧的。这就证明了线性映射 \(u\) 是紧的。

推论 1（绍德尔）. --- 设 E 是半赋范空间，F 是巴拿赫空间，E' 和 F' 是它们各自的强对偶，u 是从 E 到 F 的连续线性映射。下列性质等价：

\begin{enumerate}
\def\labelenumi{(\roman{enumi})}
\item
  从 E 到 F 的线性映射 u 是紧的；
\item
  线性映射 \({}^{t}u\)（从 \(\mathrm{F}'\) 到 \(\mathrm{E}'\) 中）是紧的；
\item
  线性映射 \({}^{t}u\)（从 \(\mathrm{F}_{c}^{\prime}\) 到 \(\mathrm{E}'\) 中）是连续的。
\end{enumerate}

(i) 与 (iii) 的等价性由命题 9 的 a) 和 b) 得出；蕴含关系 (iii)⇒(ii) 来自从 F' 到 F'c 的恒等映射是紧的这一事实（III, p. 7, 命题 7 之推论）。

最后我们来证明条件 (ii) 蕴含 (i)。假设线性映射 \({}^{t}u\)：\(F^{\prime} \to E^{\prime}\) 是紧的。把蕴含关系 (i)⇒(ii) 应用于 \({}^{t}u\)，便知 \({}^{tt}u\) 是从 \(E''\) 到 \(F''\) 的紧线性映射。用 \(v\) 表示从 E 到 \(E''\) 的典范线性映射；它是连续的。由于 F 等同于 \(F''\) 的一个闭向量子空间，且 u 在 E 上与 \({}^{t}({}^{t}u) \circ v\) 重合，故由 III, p. 3, 注记 5 推出线性映射 \(u\) 是紧的。

推论 2. --- 设 E 是半赋范空间，F 是可数型的巴拿赫空间。设 u 是从 E 到 F 的连续线性映射。假设对 F' 中元素的每个弱趋于 0 的序列 \((y_{n})\)，序列 \((^{t}u(y_{n}))\) 都在 E' 中强趋于 0。则线性映射 u 是紧的。

用 B′ 表示 F′ 的单位球，赋有由拓扑 \(\sigma(\mathrm{F}',\mathrm{F})\) 诱导的拓扑。由于 F 是可数型的巴拿赫空间，B' 是紧的可度量化空间（EVT, III, p. 19, cor. 2）。在所作假设之下，\({}^{t}u\) 在 B' 上的限制是从 B' 到 E' 的连续映射，且 \({}^{t}u(\mathrm{B}')\) 是 E' 的一个紧子集。由推论 1，线性映射 u 是紧的。

注记. --- 设 E 和 F 是分离的局部凸空间。从 E 到 F 的紧线性映射的转置并不总是从 \(F_{b}^{\prime}\) 到 \(E_{b}^{\prime}\) 的紧线性映射（参见 III, p. 108, 习题 15）。

